% !TeX root = ../../Book2.tex
% 纲要标题：### 3.4 图追踪
% 纲要位置：计划纲要.md，第 905 行。

\section{图追踪}
\label{b2:sec:0304}

在\cref{b2:lem:short-five}的证明中，我们先把元素送到相邻的模，再用正合性把一个像为零的元素写成前一映射的像。这种沿交换图施加映射、寻找原像并修正差值的证明方法，称为\term[tu-zhuizong]{图追踪}[diagram chase]。每一次提升都须说明为何存在，不同选择也须给出相同结论。本节先比较较长正合列中的同态，再构造把核与余核连接起来的映射。

% 纲要内容（仅作写作计划，尚非教材正文）：
%
% * 五引理与蛇引理的陈述
% * 连接同态的提升障碍与蛇引理的证明
% * 蛇引理的自然性与符号（3.4.3*）
%
% ---
%

\subsection{五引理与蛇引理的陈述}
\label{b2:subsec:0304:01}

短五引理的单射与满射证明可以分别放进更长的正合列。左右端不必都是零模，但相应的单射、满射条件必须补上。

\begin{theorem}[五引理]\label{b2:thm:five-lemma}
设 \(R\) 是环，下列左 \(R\) 模同态构成的交换图的两行正合：
\[
\begin{tikzcd}[column sep=1.7em,row sep=2em]
A_1\arrow[r,"a_1"]\arrow[d,"f_1"']&A_2\arrow[r,"a_2"]\arrow[d,"f_2"]&A_3\arrow[r,"a_3"]\arrow[d,"f_3"]&A_4\arrow[r,"a_4"]\arrow[d,"f_4"]&A_5\arrow[d,"f_5"]\\
B_1\arrow[r,"b_1"']&B_2\arrow[r,"b_2"']&B_3\arrow[r,"b_3"']&B_4\arrow[r,"b_4"']&B_5.
\end{tikzcd}
\]
若 \(f_1\) 满射，\(f_2,f_4\) 是同构，\(f_5\) 单射，则 \(f_3\) 是同构。
\end{theorem}

\begin{proof}
先证 \(f_3\) 单射。设 \(x\in A_3\) 且 \(f_3(x)=0\)。我们要找到 \(z\in A_1\)，把 \(x\) 写成 \(a_2a_1(z)\)，因为上行正合性给出的 \(a_2a_1=0\) 随即使 \(x=0\)。为此，先用右侧方块检查 \(x\) 能否来自 \(A_2\)：交换性给出 \(f_4(a_3(x))=b_3(f_3(x))=0\)，由 \(f_4\) 单射得 \(a_3(x)=0\)。上行在 \(A_3\) 处正合，故存在 \(y\in A_2\)，使 \(a_2(y)=x\)。于是
\[
b_2(f_2(y))=f_3(a_2(y))=0.
\]
现在要证明 \(y\) 本身来自 \(A_1\)。已知 \(f_2(y)\) 被 \(b_2\) 送到零，下行在 \(B_2\) 处正合，故可取 \(z'\in B_1\)，使 \(b_1(z')=f_2(y)\)。由 \(f_1\) 满射，可以把这个下行原像提升为 \(z\in A_1\)，使 \(f_1(z)=z'\)。交换性使 \(y\) 与 \(a_1(z)\) 在 \(f_2\) 下有相同的像，即
\[
f_2\bigl(y-a_1(z)\bigr)=f_2(y)-b_1(f_1(z))=0.
\]
由 \(f_2\) 单射，\(y=a_1(z)\)，从而 \(x=a_2a_1(z)=0\)。

再证 \(f_3\) 满射。给定 \(x'\in B_3\)，先寻找 \(x\in A_3\)，使 \(f_3(x)\) 与 \(x'\) 在 \(b_3\) 下有相同的像。这样差值就来自 \(B_2\)，而 \(f_2\) 的满射性将允许我们从 \(A_2\) 修正 \(x\)。先由 \(f_4\) 满射取 \(w\in A_4\)，使 \(f_4(w)=b_3(x')\)。要使 \(w\) 来自 \(A_3\)，需检查 \(a_4(w)=0\)；交换性给出
\[
f_5(a_4(w))=b_4(f_4(w))=b_4b_3(x')=0.
\]
\(f_5\) 单射，故 \(a_4(w)=0\)；上行在 \(A_4\) 处正合，故可取 \(x\in A_3\)，使 \(a_3(x)=w\)。于是
\[
b_3\bigl(x'-f_3(x)\bigr)=b_3(x')-f_4(a_3(x))=0.
\]
下行在 \(B_3\) 处正合，所以 \(x'-f_3(x)=b_2(y')\) 对某个 \(y'\in B_2\) 成立。由 \(f_2\) 满射取 \(y\in A_2\)，使 \(f_2(y)=y'\)。修正 \(x\) 后得到
\[
f_3\bigl(x+a_2(y)\bigr)=f_3(x)+b_2(f_2(y))=x'.
\]
因此 \(f_3\) 满射，结合已证单射性即得同构。
\end{proof}

这个结果称为\term[wu-yinli]{五引理}[five lemma]。两个方向使用的条件可以分开记录：单射证明只需 \(f_1\) 满射、\(f_2,f_4\) 单射；满射证明只需 \(f_2,f_4\) 满射、\(f_5\) 单射。它们都利用了正合性来修正一次选取，而没有要求给原像规定某种唯一的选择。

\ProseParagraphBreak
下面不再假定相邻竖直映射都是同构。若一个元素在右侧竖直映射下为零，它能否来自中间竖直映射的核？中间项中存在一个原像，并不保证这个原像已经在核中；需要测出它的竖直像，再判断能否用左侧元素把这个像修正为零。

\ProseParagraphBreak
考虑下列交换图，两行在各个中间位置正合：
\begin{equation}\label{b2:eq:snake-diagram}
\begin{tikzcd}[column sep=1.7em,row sep=2em]
{}&A\arrow[r,"u"]\arrow[d,"\alpha"']&B\arrow[r,"v"]\arrow[d,"\beta"]&C\arrow[r]\arrow[d,"\gamma"]&0\\
0\arrow[r]&A'\arrow[r,"u'"']&B'\arrow[r,"v'"']&C'&{}.
\end{tikzcd}
\end{equation}
于是 \(v\) 满射，\(u'\) 单射，\(\operatorname{im}u=\ker v\)、\(\operatorname{im}u'=\ker v'\)。这里没有要求 \(u\) 单射或 \(v'\) 满射。交换性 \(\beta u=u'\alpha\)、\(\gamma v=v'\beta\) 使 \(u,v\) 能限制在相应的核上，也使 \(u',v'\) 能下降到相应的余核上。记这些诱导映射为
\[
u_0:\ker\alpha\longrightarrow\ker\beta,\quad
v_0:\ker\beta\longrightarrow\ker\gamma,
\]
\[
\overline{u'}:\operatorname{coker}\alpha\longrightarrow\operatorname{coker}\beta,
\quad
\overline{v'}:\operatorname{coker}\beta\longrightarrow\operatorname{coker}\gamma.
\]
例如 \(\overline{u'}([a'])=[u'(a')]\)，这里两边的方括号分别取模 \(\operatorname{im}\alpha\) 与 \(\operatorname{im}\beta\) 的陪集；更换 \(a'\) 的代表元所带来的变化，经 \(u'\) 后属于 \(\operatorname{im}\beta\)，所以这个公式良定义。同样，由交换性可得
\[
v'(\operatorname{im}\beta)=v'\beta(B)=\gamma v(B)
\subseteq\operatorname{im}\gamma.
\]
因此 \([b']\mapsto[v'(b')]\) 与代表元的选择无关，确实定义了 \(\overline{v'}:\operatorname{coker}\beta\to\operatorname{coker}\gamma\)。

\ProseParagraphBreak
具体地，对 \(c\in\ker\gamma\)，先取 \(v(b)=c\)。若 \(\beta(b)\ne0\)，这一个 \(b\) 还不是所需的核中的提升。不过 \(v'(\beta(b))=\gamma(c)=0\)，所以下行正合性使 \(\beta(b)=u'(a')\)。若 \(a'=\alpha(a)\)，就可把 \(b\) 换成 \(b-u(a)\)：它在 \(v\) 下的像仍为 \(c\)，在 \(\beta\) 下的像却已变为零。因此，\(a'\) 模去 \(\operatorname{im}\alpha\) 后的类，测量的正是能否把 \(c\) 提升到 \(\ker\beta\) 的障碍。

\begin{theorem}[蛇引理]\label{b2:thm:snake-lemma}
设 \(R\) 是环，\(A\xrightarrow{u}B\xrightarrow{v}C\to0\) 与 \(0\to A'\xrightarrow{u'}B'\xrightarrow{v'}C'\) 是左 \(R\) 模的正合列，同态 \(\alpha:A\to A'\)、\(\beta:B\to B'\)、\(\gamma:C\to C'\) 满足 \(\beta u=u'\alpha\)、\(\gamma v=v'\beta\)。记 \(u_0,v_0\) 为 \(u,v\) 在竖直映射的核上诱导的同态，\(\overline{u'},\overline{v'}\) 为 \(u',v'\) 在竖直映射的余核上诱导的同态。此时存在同态
\(\delta:\ker\gamma\to\operatorname{coker}\alpha\)，使下列序列正合：
\begin{equation}\label{b2:eq:snake-exact-sequence}
\begin{tikzcd}[column sep=1.2em]
\ker\alpha\arrow[r,"u_0"]&\ker\beta\arrow[r,"v_0"]&\ker\gamma\arrow[r,"\delta"]&\operatorname{coker}\alpha\arrow[r,"\overline{u'}"]&\operatorname{coker}\beta\arrow[r,"\overline{v'}"]&\operatorname{coker}\gamma.
\end{tikzcd}
\end{equation}
对 \(c\in\ker\gamma\)，由 \(v\) 满射选 \(b\in B\) 使 \(v(b)=c\)。此时 \(\beta(b)\in\ker v'=\operatorname{im}u'\)，下行正合性保证存在 \(a'\in A'\) 满足 \(u'(a')=\beta(b)\)，而 \(u'\) 单射保证这个 \(a'\) 对给定的 \(b\) 唯一。连接同态由
\begin{equation}\label{b2:eq:snake-connecting-definition}
\delta(c)=a'+\operatorname{im}\alpha
\end{equation}
给出，所得余核类不依赖于提升 \(b\) 的选择。
若 \(u\) 还单射，可在\eqref{b2:eq:snake-exact-sequence}的左端添 \(0\)；若 \(v'\) 还满射，可在右端添 \(0\)。特别地，两行均为短正合列时，两端都可添 \(0\)。
\end{theorem}

这一结果称为\term[she-yinli]{蛇引理}[snake lemma]，其中的 \(\delta\) 称为\term[lianjie-tongtai]{连接同态}[connecting homomorphism]。构造时从 \(c\) 出发，沿 \(v\) 反向选一个原像 \(b\)，施加 \(\beta\)，再沿 \(u'\) 反向取原像 \(a'\)，最后取余核类：
\[
\begin{tikzcd}[column sep=2.5em,row sep=1.5em]
c&b\arrow[l,"v"']\arrow[d,"\beta"]&&\\
&\beta(b)&a'\arrow[l,"u'"]\arrow[r,mapsto]&{[a']}.
\end{tikzcd}
\]
图中的箭头仍按原同态的方向画出；反向寻找原像并不是给出逆同态，因而这个公式不能视为原图中几个箭头的直接复合。

\subsection{连接同态与蛇引理的证明}
\label{b2:subsec:0304:02}

\begin{proof}[\cref{b2:thm:snake-lemma}的证明]
先构造 \(\delta\)。给定 \(c\in\ker\gamma\)，由 \(v\) 满射选 \(b\) 使 \(v(b)=c\)。交换性给出
\[
v'(\beta(b))=\gamma(v(b))=\gamma(c)=0.
\]
由下行在 \(B'\) 处正合，\(\beta(b)\in\operatorname{im}u'\)；由 \(u'\) 单射，存在唯一的 \(a'\in A'\) 满足 \(u'(a')=\beta(b)\)。按\eqref{b2:eq:snake-connecting-definition}定义 \(\delta(c)\)。

要证明 \(\delta(c)\) 与提升的选择无关，就需证明另一个提升所得的 \(a'_1-a'\) 属于 \(\operatorname{im}\alpha\)。若另选 \(b_1\) 满足 \(v(b_1)=c\)，则 \(b_1-b\in\ker v=\operatorname{im}u\)，所以 \(b_1-b=u(a)\) 对某个 \(a\in A\) 成立。若 \(a'_1\) 是由 \(b_1\) 得到的元素，则
\[
u'(a'_1-a')=\beta(b_1-b)=\beta(u(a))=u'(\alpha(a)).
\]
\(u'\) 单射，故 \(a'_1-a'=\alpha(a)\)。两者在 \(\operatorname{coker}\alpha\) 中给出同一陪集。这里 \(a'\) 本身可能随提升改变，保持不变的是它模去 \(\operatorname{im}\alpha\) 后的类。

检验 \(R\) 线性。对 \(c_1,c_2\in\ker\gamma\)，分别选提升 \(b_1,b_2\) 与相应元素 \(a'_1,a'_2\)。此时 \(b_1+b_2\) 是 \(c_1+c_2\) 的提升，\(a'_1+a'_2\) 满足所需的 \(u'\) 像等式，所以 \(\delta(c_1+c_2)=\delta(c_1)+\delta(c_2)\)。对 \(x\in R\)，取 \(xb_1\) 与 \(xa'_1\)，得到 \(\delta(xc_1)=x\cdot\delta(c_1)\)。这些计算只使用左模同态的线性性，不移动标量的次序。

在 \(\ker\beta\) 处，要证 \(\operatorname{im}u_0=\ker v_0\)。由 \(vu=0\)，有 \(\operatorname{im}u_0\subseteq\ker v_0\)。反过来，若 \(b\in\ker\beta\) 且 \(v(b)=0\)，上行正合使 \(b=u(a)\)。于是 \(u'(\alpha(a))=\beta(u(a))=0\)，由 \(u'\) 单射得 \(\alpha(a)=0\)，故 \(b\in\operatorname{im}u_0\)。

在 \(\ker\gamma\) 处，要证 \(\operatorname{im}v_0=\ker\delta\)。若 \(c=v_0(b)\)，其中 \(\beta(b)=0\)，则构造中可取 \(a'=0\)，所以 \(\delta(c)=0\)。反过来，若 \(\delta(c)=0\)，选 \(b,a'\) 如定义。此时 \(a'\in\operatorname{im}\alpha\)，可写 \(a'=\alpha(a)\)。修正提升为 \(b-u(a)\)，就有
\[
\beta\bigl(b-u(a)\bigr)=u'(a')-u'(\alpha(a))=0,
\qquad v\bigl(b-u(a)\bigr)=c.
\]
故 \(c\in\operatorname{im}v_0\)，证明 \(\delta(c)=0\) 当且仅当这种核中的提升存在。

在 \(\operatorname{coker}\alpha\) 处，要证 \(\operatorname{im}\delta=\ker\overline{u'}\)。若 \(\delta(c)=[a']\)，构造式 \(u'(a')=\beta(b)\) 表明 \(\overline{u'}([a'])=0\)。反过来，若 \(\overline{u'}([a'])=0\)，则存在 \(b\in B\)，使 \(u'(a')=\beta(b)\)。令 \(c=v(b)\)。因为
\[
\gamma(c)=\gamma(v(b))=v'(\beta(b))=v'(u'(a'))=0,
\]
有 \(c\in\ker\gamma\)，而选取这一个 \(b\) 计算连接同态，就得 \(\delta(c)=[a']\)。

在 \(\operatorname{coker}\beta\) 处，要证 \(\operatorname{im}\overline{u'}=\ker\overline{v'}\)。由 \(v'u'=0\)，有 \(\operatorname{im}\overline{u'}\subseteq\ker\overline{v'}\)。若 \(\overline{v'}([b'])=0\)，就有 \(v'(b')=\gamma(c)\) 对某个 \(c\in C\) 成立。用 \(v\) 满射选 \(b\in B\)，使 \(v(b)=c\)。于是
\[
v'\bigl(b'-\beta(b)\bigr)=v'(b')-\gamma(v(b))=0.
\]
下行正合，所以 \(b'-\beta(b)=u'(a')\) 对某个 \(a'\in A'\) 成立。模去 \(\operatorname{im}\beta\) 后，得到 \([b']=\overline{u'}([a'])\)，证明反向包含。

最后，若 \(u\) 单射，其限制 \(u_0\) 也单射，故左端可以添零。若 \(v'\) 满射，对 \(\operatorname{coker}\gamma\) 中任意代表元 \(c'\in C'\)，选 \(b'\in B'\) 满足 \(v'(b')=c'\)，则 \(\overline{v'}([b'])=[c']\)，故右端可以添零。
\end{proof}

连接同态可以非零。取整数 \(n\ge2\)，将\eqref{b2:eq:integer-short-exact}在上下两行各写一次，取竖直映射
\[
\alpha:\mathbb Z\xrightarrow{\times n}\mathbb Z,
\qquad\beta:\mathbb Z\xrightarrow{\times n}\mathbb Z,
\qquad\gamma:\mathbb Z/n\mathbb Z\xrightarrow{0}\mathbb Z/n\mathbb Z.
\]
左方块交换，因为两个方向都乘以 \(n^2\)；右方块交换，因为整数的 \(n\) 倍在模 \(n\) 后为零。此时 \(\ker\gamma=\mathbb Z/n\mathbb Z\)，\(\operatorname{coker}\alpha=\mathbb Z/n\mathbb Z\)。对 \(c=a+n\mathbb Z\)，选上行中间项的提升 \(b=a\)，则 \(\beta(b)=na=u'(a)\)，所以
\[
\delta(a+n\mathbb Z)=a+n\mathbb Z.
\]
连接同态正是恒等同构。这里 \(\ker\beta=0\)，而 \(\ker\gamma=\mathbb Z/n\mathbb Z\ne0\)，所以每个非零的 \(c\in\ker\gamma\) 都没有来自 \(\ker\beta\) 的提升，并且 \(\delta(c)=c\ne0\)。连接同态给出的同构 \(\ker\gamma/\operatorname{im}v_0\cong\operatorname{im}\delta\)，在这个例子中就是把整个右侧核送到左侧余核。

\OptionalSubsection{蛇引理的自然性与符号}{b2:subsec:0304:03}

现在比较两个满足蛇引理条件的图：第二个图的对象和映射用波浪号表示，用 \(s_\bullet\) 记从第一个图上排到第二个图上排的同态，用 \(t_\bullet\) 记下排到下排的同态。它们须与水平和竖直映射都交换，才构成两个图之间的相容映射。连接同态的构造不依赖选定一套原像。因而，先取第一个连接同态再经下排映射送出，或先经上排映射送出再取第二个连接同态，应该在第二个图的余核中得到同一个类。下面将这个相容性写成具体等式。

\begin{proposition}[蛇引理的自然性]\label{b2:prop:snake-naturality}
设 \(R\) 是环，给定两个左 \(R\) 模的交换图，每个图的上行形如 \(A\xrightarrow{u}B\xrightarrow{v}C\to0\)，下行形如 \(0\to A'\xrightarrow{u'}B'\xrightarrow{v'}C'\)，两行正合，竖直映射 \(\alpha,\beta,\gamma\) 与水平映射交换；第二个图的对象和映射用波浪号表示。设有六个同态
\[
s_A:A\to\widetilde A,\quad s_B:B\to\widetilde B,\quad
s_C:C\to\widetilde C,
\]
\[
t_A:A'\to\widetilde A',\quad t_B:B'\to\widetilde B',\quad
t_C:C'\to\widetilde C',
\]
它们与两个图中的所有箭头相容，即
\[
\begin{aligned}
s_Bu&=\widetilde u s_A,&s_Cv&=\widetilde v s_B,\\
t_Bu'&=\widetilde u' t_A,&t_Cv'&=\widetilde v' t_B,
\end{aligned}
\]
以及 \(t_A\alpha=\widetilde\alpha s_A\)、\(t_B\beta=\widetilde\beta s_B\)、\(t_C\gamma=\widetilde\gamma s_C\)。此时它们诱导两条蛇正合列之间的相容同态。特别地，若 \(s_{C,0}\) 是 \(s_C\) 在 \(\ker\gamma\) 上的限制，\(\overline{t_A}\) 是 \(t_A\) 诱导的余核映射，则
\begin{equation}\label{b2:eq:snake-naturality}
\overline{t_A}\,\delta=\widetilde\delta\,s_{C,0}.
\end{equation}
\end{proposition}

\begin{proof}
若 \(c\in\ker\gamma\)，有 \(\widetilde\gamma(s_C(c))=t_C(\gamma(c))=0\)，所以 \(s_{C,0}\) 的值确实位于 \(\ker\widetilde\gamma\)。另一方面，\(t_A\alpha=\widetilde\alpha s_A\) 保证 \(t_A\) 把 \(\operatorname{im}\alpha\) 送进 \(\operatorname{im}\widetilde\alpha\)，因而 \(\overline{t_A}([a'])=[t_A(a')]\) 良定义。其他核与余核之间的映射以同样的限制或商映射公式定义；与 \(u_0,v_0,\overline{u'},\overline{v'}\) 的相容性直接由所列水平等式限制或下降得到。

只需核验含连接同态的方块。对 \(c\in\ker\gamma\)，选择 \(b,a'\) 满足 \(v(b)=c\)、\(u'(a')=\beta(b)\)。在第二个图中，\(s_B(b)\) 是 \(s_C(c)\) 的一个提升，因为
\[
\widetilde v(s_B(b))=s_C(v(b))=s_C(c).
\]
又有
\[
\widetilde\beta(s_B(b))
=t_B(\beta(b))
=t_B(u'(a'))
=\widetilde u'(t_A(a')).
\]
故使用这组允许的提升计算 \(\widetilde\delta\)，得到
\[
\widetilde\delta(s_C(c))=[t_A(a')]
=\overline{t_A}([a'])=\overline{t_A}(\delta(c)).
\]
这就是\eqref{b2:eq:snake-naturality}。提升的选择无关性保证我们可以使用这一组特别便于比较的提升。
\end{proof}

式\eqref{b2:eq:snake-naturality}把连接同态也放入两条蛇正合列之间的交换方块：
\[
\begin{tikzcd}[column sep=4em,row sep=2.5em]
\ker\gamma\arrow[r,"\delta"]\arrow[d,"s_{C,0}"']&\operatorname{coker}\alpha\arrow[d,"\overline{t_A}"]\\
\ker\widetilde\gamma\arrow[r,"\widetilde\delta"']&\operatorname{coker}\widetilde\alpha.
\end{tikzcd}
\]
自然性在这里有明确的计算作用：比较两个问题时，不必为两边各自任意选择提升，再设法比较；可以把一边已经选好的提升送到另一边去。定义的选择无关性保证这样得到的仍是原先那个连接同态。

\ProseParagraphBreak
本书在\eqref{b2:eq:snake-connecting-definition}中固定正号：满足 \(u'(a')=\beta(b)\) 时，令 \(\delta(c)=[a']\)。如果把连接同态统一改定义为 \(-\delta\)，其核与像都不变，所以蛇正合列仍然正合，两个图同时采用负号时自然性等式也仍成立。但是比较具体公式时必须采用同一约定。例如将下行首映射 \(u'\) 换成 \(-u'\)，并把竖直映射 \(\alpha\) 同时换成 \(-\alpha\)，才能保持左方块交换；在原余核的同一陪集记号下，新连接同态就是 \(-\delta\)，因为新的原像为 \(-a'\)。

\ProseParagraphBreak
因此，使用图追踪所得的同态时，应保留定义它的交换图、箭头方向与元素公式。正合性提供存在性，单射性提供需要的唯一性，商映射消去选择造成的差别；这些条件各有具体的使用位置。后面遇到更长的正合列时，仍会反复使用这几步。
